\scanpage{229}
\mspage{179}
% L1
(\unsure{p. ex|op. cit|loc. cit}) $\pi_i(Z)=0$ pour $i\in[0,n]$.

% L2–5
$b_n$) \inserted{$\pi_0(f)$ est bijectif, et} $\forall\,x\deleted{\unsure{\in\mathrm{Ob}}}\unsure{\in}X$, $\pi_i(f,x):\pi_i(X,x)\unsure{-|\to}\pi_i(Y,f(x))$ est iso pour $i\in[\unsure{1|0},n]$. (\uline{NB} Il suffit, \deleted{\unsure{pour le dire à ce}} de le dire pour \uline{un} point $x$ de chaque composante connexe de $X$.)

% L6–8
$c_n$) $H^i(Y,F)\longrightarrow H^i(X,f^*(F))$ iso pour $F$ fais. loc. constant sur $Y$ (avec $F$ abélien si $i\geq2$, $F$ faisceau en groupes si $i=1$).

% L9
On a alors
% L10
\[
a_n\Longrightarrow b_n\deleted{\unsure{b}}\overset{\inserted{\text{\parbox{5.6cm}{si $X,Y$ 1-connexes c'est vraiment prouvé, sinon, il y a des vérifications à faire (\unsure{qui ne|qui ont} \unsure{l'air|ainsi} \unsure{soit elles|soit des}\ldots)}}}}{\Longrightarrow}c_n\overset{*}{\Longrightarrow}a_{n-1}
\]
% L10, left margin
% slanted
\margin{$^*$\uline{NB} Même si $X,Y$ 1-connexes, il n'est pas clair que $c_n\Rightarrow b_n$, car il faut montrer que $\pi_n(Z)\to\pi_n(X)$ est nul, et on trouve seulement que le composé avec $\pi_n(X)\to H_n(X)$ est nul. Je présume qu'il y a des contre-exemples.}
% L11–12
(et je présume que chacune de ces implications est stricte).

% L13–16
\uline{Dém} $a_n\Rightarrow b_n$ est conséquence immédiate de la suite exacte d'homotopie, pour une $\mathrm{Hot}$-fibration $\overline X\xrightarrow{\bar f}Y$ qui factorise $X\xrightarrow{f}Y$ en $f=\bar f\circ i$, $i\in W_\infty$. En fait,
% L17–22
\deleted{\unsure{$b_n$}}\inserted{$a_n$} \uline{équivaut} à $b_n$ $+$ la surjectivité des $\pi_{n+1}(X,x)\to\pi_{n+1}(Y,f(x))$, laquelle bien sûr ne résulte \uline{pas} de $b_n$, i.e. de la bijectivité de $\pi_i(f,x)$ pour $i\leq n$ (prendre des $K(\pi,n+1)$). Donc l'implication est bien stricte.
% archive sheet number: 229
